% Generated draft: requires mathematical and source review.
\documentclass[11pt]{article}
\usepackage[margin=25mm]{geometry}
\usepackage{fontspec}
\setmainfont{Latin Modern Roman}
\usepackage{amsmath,amssymb,mathtools,unicode-math}
\setmathfont{Latin Modern Math}
\usepackage{enumitem,needspace,xurl,hyperref}
\setlength{\parindent}{0pt}
\setlength{\parskip}{5pt}
\setlength{\emergencystretch}{3em}
\newcommand{\N}{\mathbb{N}}
\newcommand{\Z}{\mathbb{Z}}
\newcommand{\Q}{\mathbb{Q}}
\newcommand{\R}{\mathbb{R}}
\newcommand{\C}{\mathbb{C}}
\newcommand{\F}{\mathbb{F}}
\newcommand{\A}{\mathbb{A}}
\newcommand{\PP}{\mathbb{P}}
\newcommand{\E}{\mathbb{E}}
\newcommand{\eps}{\varepsilon}
\newcommand{\dd}{\,\mathrm{d}}
\begin{document}
\section*{UnsolvedMath ID: 800041}
\textbf{Source problem code:} catalogue entry
\subsection{Definitions and mathematical statement}
A semigroup is a nonempty set equipped with an associative binary operation; this uses the usual convention, since the supplied record does not specify whether empty semigroups are allowed. An isomorphism of semigroups is a bijection preserving multiplication.

The supplied problem asks for semigroups in the ``basic classes of semigroups'' that are determined, up to isomorphism, by their ``identical inclusions.'' Neither quoted expression is defined in the supplied record. Consequently, the following is a parameterized formulation of the requested determination property, pending specification of those expressions and of the comparison class.

Let \(\mathcal{B}\) be a specified basic class of semigroups, let \(\mathcal{C}\) be a specified comparison class containing \(\mathcal{B}\), and, once the syntax and satisfaction relation for identical inclusions have been specified, write
\[
\mathcal{I}(S)=\{\varphi : \varphi\text{ is an identical inclusion satisfied by }S\}.
\]
Determine all \(S\in\mathcal{B}\) such that
\[
\forall T\in\mathcal{C},\qquad
\mathcal{I}(T)=\mathcal{I}(S)\ \Longrightarrow\ T\cong S.
\]
Here equality means satisfaction of exactly the same identical inclusions. The record does not specify whether comparison is with all semigroups or only with members of a particular class, nor does it identify the basic classes to be considered. Thus this schema does not resolve the missing scope or the undefined notion of identical inclusion.

As background, the supplied source reports that the two-element semilattice (reference [L1]), quasicyclic groups (reference [BM]), and every finite group (reference [M2]) are determined by identical inclusions up to isomorphism. These are reported source claims, whose precise hypotheses and comparison classes cannot be recovered from the supplied record. The registry labels the problem open; that label is supplied metadata, not an independently verified mathematical assessment.
\subsection{Short English statement}
Identify semigroups in the source's basic classes that are uniquely determined up to isomorphism by all identical inclusions they satisfy. A complete formulation requires the definition of identical inclusion, the intended basic classes, and the class of semigroups against which uniqueness is tested. The source reports the two-element semilattice, quasicyclic groups, and all finite groups as known examples.
\subsection{Sources}
\begin{itemize}
\item Local UnsolvedMath registry, permanent ID \(800041\), exact source problem code \texttt{catalogue entry}, title ``The Lyapin's Notebook --- Problem 43.'' The supplied record labels the problem open and lists Bershadsky, Kublanovsky, and Mashevitzky under proposed\_by, with proposed year 2026.
\item Bershadsky, Kublanovsky, Mashevitzky, \emph{The Lyapin's Notebook: Unsolved Problems in Semigroup Theory} (2026), Problem 43, attributed in the supplied statement to G. Mashevitzky. Supplied source URL: \url{https://arxiv.org/abs/2604.04763}. The record gives access date 2026-10-16 and extraction method \texttt{arxiv\_latex}; these are supplied provenance, not a claim of independent access or verification.
\item Reference [L1] is supplied for the determination of the two-element semilattice; reference [BM] for quasicyclic groups; reference [M2] for every finite group. Their bibliographic entries and precise results are absent from the supplied record. The notebook is the supplied external citation; the underlying references have not been independently identified or verified.
\end{itemize}
\end{document}
